\section{Testing a trimmed null, and an analytic baseline}
\label{sec:analytic}

For each level $\alpha\in[0,1)$ we test
\begin{equation}
H_{0,\alpha}:\ \TV(P,Q)\le\alpha
\qquad\text{against}\qquad
H_{1,\alpha}:\ \TV(P,Q)>\alpha ,
\label{eq:null}
\end{equation}
and by \eqref{eq:zeroeq} the null holds exactly when $D_\alpha(P,Q)=0$. The null is a total-variation neighborhood of
$P=Q$, a relevant hypothesis in the sense of \citet{dettekokotaue2020}, as in the trimming-based tests of
\citet{alvarez2012similarity} and \citet{delbarrio2020ks} and in robust tests between neighborhoods of laws
\citep{huber1965robust,huber1973minimax,schrab2025robust,liubriol2025robust}, which fix its size in advance, whereas
here it is the quantity to be estimated. Kernel tests of composite nulls \citep{key2025composite} calibrate over a
parametric family, and equivalence tests \citep{wellek2010equivalence} put such a neighborhood in the alternative
(Section~\ref{app:analytic} of the supplement). This section gives a threshold $\tau(\alpha)$, computed before the
data, that $\widehat D_\alpha$ exceeds with probability at most $\delta$ under every pair in the null. Inverted over
$\alpha$, it gives a baseline that needs no split but far more observations than our method.

\subsection{A threshold for the whole null}

Limit laws are of little help in calibrating $\widehat D_\alpha$, because at the boundary of the null the two regions
touch, and a minimum over trimmings is then only directionally differentiable in the laws, with a limit law that
depends on the unknown geometry of the contact \citep{fang2019inference}, and permutation is exact only under $P=Q$. We
bound $\widehat D_\alpha$ instead by a difference of sample means. Under $H_{0,\alpha}$ the laws share an
$\alpha$-trimming $R$ by \eqref{eq:tvintersect}, with densities $g_P=dR/dP$ and $g_Q=dR/dQ$ bounded by $\beta$, so the
weights $w_i=g_P(X_i)/n$ and $v_j=g_Q(Y_j)/m$ respect the caps of \eqref{eq:qp}, and the embeddings they give are
averages of independent vectors with the same mean $\mu_R$. Their totals fluctuate around one, and rescaling the weights,
or adding mass under the caps, gives a feasible pair at the price of a term proportional to the shortfalls
(Lemma~\ref{lem:plugin} of the supplement). So $\widehat D_\alpha$ is at most the norm of a centered difference of two
averages of vectors bounded by $\beta\kappa$, plus two shortfalls, and bounded-difference and Bernstein inequalities
give a threshold in closed form. Let $n_\star=nm/(n+m)$, so that $1/n_\star=1/n+1/m$, $c=\max(1/n,1/m)$, $L=\log(6/\delta)$ and $\bar\kappa^2=\min(\kappa^2,\Delta_\phi^2/2)$.

\begin{theorem}[Plug-in threshold]
\label{thm:plugin}
Let (A1) and (A2) of Assumption~\ref{ass:regularity} hold, and let
\begin{gather*}
\tau^{\mathrm P}(0)=\bar\kappa/\sqrt{n_\star}+\Delta_\phi\sqrt{\log(1/\delta)/(2n_\star)},\\
\tau^{\mathrm P}(\alpha)=\sqrt{\beta \bar\kappa^2/n_\star}+\beta \Delta_\phi\sqrt{L/(2n_\star)}+\Delta_\phi\Big[\sqrt{2(\beta-1)L/n_\star}+\frac{cL}3\Big]
\qquad(0<\alpha<1).
\end{gather*}
Then for every $\alpha\in[0,1)$, $n,m\ge1$, $\delta\in(0,1)$ and every pair with $\TV(P,Q)\le\alpha$,
$\Pp\big(\widehat D_\alpha>\tau^{\mathrm P}(\alpha)\big)\le\delta$. Moreover, $\alpha\mapsto\tau^{\mathrm P}(\alpha)$
is nondecreasing on $[0,1)$.
\end{theorem}

The threshold depends on $\kappa$, $\Delta_\phi$, $\alpha$, $n$, $m$ and $\delta$ alone and is, up to a term of order
$\Delta_\phi\log(1/\delta)/(n\wedge m)$, of order $\beta \Delta_\phi\{(1+\log(1/\delta))/(n\wedge m)\}^{1/2}$. It
jumps up at $0+$, where $L>\log(1/\delta)$ replaces $\log(1/\delta)$. The pair enters only through the existence of $R$. At
$\alpha=0$ the common trimming is $R=P=Q$, and Theorem~\ref{thm:plugin} is a finite-sample bound on the MMD of the
kind given by \citet[Corollary~9]{gretton2012kernel}. For the Gaussian kernel and $n=m$ it reads
$\tau^{\mathrm P}(0)=[1+\sqrt{\log(1/\delta)}\,]\sqrt{2/n}$.

We write $\tau=\tau^{\mathrm P}$ for the threshold of the test. A second valid threshold $\tau^{\mathrm S}(\alpha,\delta)$
comes from the two-sided bound of Theorem~\ref{thm:twosided} of the supplement, by the same reduction to means as
\eqref{eq:hausdorff}, but it is larger than $\tau^{\mathrm P}$ at the sample sizes we use. Since $\tau$ is
nondecreasing and the empirical path is nonincreasing, one event serves every level.

\begin{lemma}[One event for all levels]
\label{lem:oneevent}
Let $\tau:[0,1)\to[0,\infty)$ be nondecreasing, and suppose that $\Pp(\widehat D_\alpha>\tau(\alpha))\le\delta$ for
every $\alpha$ and every pair with $\TV(P,Q)\le\alpha$. Then for every pair with $\alpha^\star<1$, the event
$\{\widehat D_{\alpha^\star}\le\tau(\alpha^\star)\}$ has probability at least $1-\delta$, and on it
$\widehat D_\alpha\le\tau(\alpha)$ for every $\alpha\in[\alpha^\star,1)$.
\end{lemma}

The analytic test rejects $H_{0,\alpha}$ when
\begin{equation}
\widehat D_\alpha>\tau(\alpha) .
\label{eq:test}
\end{equation}
It has level $\delta$ for all $n,m\ge1$, uniformly over the null, by Lemma~\ref{lem:oneevent} one event carries the
guarantee for all levels at once, and its level does not use (A3). Its validity at every sample size, uniformly over the null, answers
the question of \citet{alvarez2012similarity} recalled in the introduction, as does, more simply, the betting test of
Section~\ref{sec:split}. The kernel must be fixed before the data are seen, for instance on a held-out pilot sample
(Section~\ref{sec:data}). One threshold for every pair in the null is a worst case: at the boundary of a
two-dimensional design, at $n=m=300$ and $\alpha=\alpha^\star$, $\tau(\alpha)$ is $\stCalibRatioTauP$ times the
empirical $1-\delta$ quantile of $\widehat D_\alpha$ (Section~\ref{app:experiments} of the supplement).

\subsection{The estimator}
\label{sec:alpha}

Inverting the test over $\alpha$, we define the analytic estimator, with a subscript that sets it apart from the
split-sample bound $\widehat\alpha$ of Section~\ref{sec:lower},
\begin{equation}
\widehat\alpha_{\mathrm{an}}\coloneqq\inf\big\{\alpha\in[0,1):\ \widehat D_\alpha\le\tau(\alpha)\big\} ,
\label{eq:alphaest}
\end{equation}
the first level at which the empirical path falls below the threshold. Since $\alpha\mapsto\widehat D_\alpha$ is
nonincreasing and $\tau$ nondecreasing, every level above $\widehat\alpha_{\mathrm{an}}$ is accepted and every level below it rejected,
so bisection brackets $\widehat\alpha_{\mathrm{an}}$. It is well defined, because
$\widehat D_\alpha\le \Delta_\phi$ while $\tau(\alpha)\to\infty$ as $\alpha\to1$ when $\Delta_\phi>0$. On the event of
Lemma~\ref{lem:oneevent}, of probability at least $1-\delta$ whenever $\alpha^\star<1$, every level from $\alpha^\star$
on is accepted, so $\widehat\alpha_{\mathrm{an}}\le\alpha^\star$ and every rejected level lies below $\alpha^\star$, at every sample
size, for every pair and however the levels tested were chosen. A grid of levels read upward, as in
Section~\ref{sec:data}, therefore gives a lower confidence bound too. It is consistent from below,
$\Pp(\widehat\alpha_{\mathrm{an}}<\alpha)\to0$ for $\alpha<\alpha^\star$, as $\widehat D_\alpha\ge
D_\alpha-\tau^{\mathrm S}(\alpha,\delta')$ with probability $1-\delta'$ (Theorem~\ref{thm:twosided} of the supplement)
and both thresholds tend to zero.

\subsection{How many observations}
\label{sec:samplesize}

Because $\tau$ does not depend on the data, the sample size at which the test fires can be computed before any data
are collected, given a guess at the effect. For the Gaussian kernel and $n=m$, Theorem~\ref{thm:plugin} gives
$\tau(0)\le C_0\sqrt{2/n}$, so the test rejects $H_{0,0}$ as soon as
\begin{equation}
n>2\Big(\frac{C_0}{\widehat D_0}\Big)^2,
\qquad
C_0=\bar\kappa+\Delta_\phi\sqrt{\log(1/\delta)/2}=1+\sqrt{\log(1/\delta)}\approx\cstCzero
\ \text{ at }\delta=0.05 .
\label{eq:samplesize}
\end{equation}
With a guess at $D_0=\MMD(P,Q)$ in place of $\widehat D_0$, \eqref{eq:samplesize} is a design rule, which scales like
$D_0^{-2}$.
